\documentclass[11pt,reqno]{article}

\usepackage[margin=1.08in,headheight=14pt]{geometry}
\usepackage[T1]{fontenc}
\usepackage{lmodern}
\usepackage{microtype}
\usepackage{amsmath,amssymb,amsthm,mathtools}
\usepackage{booktabs,tabularx,array}
\usepackage{fancyhdr}
\usepackage{xcolor}
\usepackage{hyperref}

\definecolor{linkblue}{RGB}{24,65,117}
\hypersetup{
  unicode=true,
  colorlinks=true,
  linkcolor=linkblue,
  citecolor=linkblue,
  urlcolor=linkblue,
  pdftitle={On Erd\H{o}s's Multiplicative Representation Problem},
  pdfauthor={Rishikesh Gajjala},
  pdfsubject={Extremal multiplicative number theory},
  pdfkeywords={multiplicative representation, multiplicative Sidon set, semiprime}
}

\linespread{1.04}
\allowdisplaybreaks
\setlength{\parindent}{1.25em}
\setlength{\parskip}{0pt}
\setlength{\emergencystretch}{2em}

\newtheorem{theorem}{Theorem}[section]
\newtheorem{proposition}[theorem]{Proposition}
\newtheorem{lemma}[theorem]{Lemma}
\newtheorem{corollary}[theorem]{Corollary}
\newtheorem{claim}[theorem]{Claim}
\theoremstyle{definition}
\newtheorem{definition}[theorem]{Definition}
\newtheorem{remark}[theorem]{Remark}

\newcommand{\N}{\mathbb N}
\newcommand{\Ucal}{\mathcal U}
\newcommand{\floor}[1]{\left\lfloor #1\right\rfloor}
\newcommand{\abs}[1]{\left|#1\right|}

\pagestyle{fancy}
\fancyhf{}
\fancyhead[L]{\small Erd\H{o}s's multiplicative representation problem}
\fancyhead[R]{\small\thepage}
\renewcommand{\headrulewidth}{0.35pt}
\fancypagestyle{plain}{
  \fancyhf{}
  \fancyhead[L]{\small Erd\H{o}s's multiplicative representation problem}
  \fancyhead[R]{\small\thepage}
  \renewcommand{\headrulewidth}{0.35pt}
}

\title{\bfseries On Erd\H{o}s's\\
Multiplicative Representation Problem}
\author{Rishikesh Gajjala\\[2pt]
\small NYU Abu Dhabi}
\date{}

\begin{document}

\maketitle
\thispagestyle{plain}

\begin{abstract}
Let \(g_3(n)\) be the largest size of a set \(A\subseteq[n]\) for which every
integer has at most two representations \(m=ab\) with \(a,b\in A\) and
\(a<b\).
We prove
\[
  g_3(n)
  =\frac{n\log\log n}{\log n}
   +\left(1+M+\Gamma+o(1)\right)\frac{n}{\log n},
\]
where \(M\) is the Meissel--Mertens constant and \(\Gamma\) is a finite
constant defined by a discrete cofactor problem.  This resolves a
five-decade-old problem posed by Erd\H{o}s. The initial proofs were generated
and formalized in Lean by OpenAI's GPT 5.6 Sol; the authors have verified and
rewritten proofs to enhance readability and provide additional context.
\end{abstract}

\medskip
\noindent\textbf{2020 Mathematics Subject Classification.}
Primary 11B83; secondary 05D05, 11N37.

\smallskip
\noindent\textbf{Keywords.}
Multiplicative representation, multiplicative Sidon set, extremal
hypergraph, semiprime, variational constant.

\section{Introduction}

Write \([n]=\{1,\ldots,n\}\), and let
\[
  \rho_A(m)=\#\{(a,b):a,b\in A,\ a<b,\ ab=m\}
\]
and
\[
  g_3(n)=\max\{|A|:A\subseteq[n],\ \rho_A(m)\leq2\text{ for all }m\}.
\]
Erd\H{o}s proved \cite{Erdos1964Multiplicative} that
\[
  g_3(n)\sim\frac{n\log\log n}{\log n}.
\]
Erd\H{o}s posed the second-order question underlying Problem 796 in his
1969 paper \emph{Problems and results in chromatic graph theory}
\cite{Erdos1969Chromatic}.  In its corrected modern form
\cite{BloomErdos796}, it asks whether
\[
  g_3(n)=\frac{n\log\log n}{\log n}
  +(c+o(1))\frac{n}{\log n}.
\]

Let
\[
  M=\lim_{x\to\infty}
  \left(\sum_{p\leq x}\frac1p-\log\log x\right)
\]
be the Meissel--Mertens constant.  Our main result is the following.

In a January 2026 note, Tang \cite{Tang2026Erdos796} identified the corrected
second-order scale and gave a construction proving
\[
  g_3(n)\geq\frac{n\log\log n}{\log n}
  +(1+M+o(1))\frac n{\log n}.
\]
We prove that the second-order constant exists, identify it as
\(1+M+\Gamma\), and give the explicit bounds
\[
  \frac4{15}\leq\Gamma<13.
\]

\begin{theorem}\label{thm:main}
There is a finite constant \(\Gamma\) satisfying
\[
  \frac4{15}\leq\Gamma<13
\]
such that
\[
  g_3(n)
  =\frac{n\log\log n}{\log n}
   +\left(1+M+\Gamma+o(1)\right)\frac{n}{\log n}.
\]
Thus the answer is yes, with \(c=1+M+\Gamma\), and
\[
  1.5281638795\ldots
  =1+M+\frac4{15}\leq c<15.
\]
\end{theorem}

The proof has three parts.  First, a pruning argument leaves two normal
forms: a large prime times a smooth cofactor, and a small multiplier times
two large primes.  A \(C_4\)-overlap argument absorbs the second form into
the semiprime baseline, while the first is encoded by compatible cofactor
families.  Finally, the resulting finite optimization problems converge to
a fixed constant \(\Gamma\).

The pruning and complete-box mechanism refine Erd\H{o}s's proof of the
leading-order asymptotic \cite[Lemma, pp.~252--253; proof of Theorem~3,
pp.~255--261]
{Erdos1964Multiplicative}.  Erd\H{o}s already split integers by the sizes of
their factors and used a multipartite hypergraph box to force many equal
products.  Here that strategy is sharpened to retain terms of order
\(n/\log n\).  The additional ingredients are the \(C_4\)-overlap
normalization, the compatible-cofactor variational model, the global
collision cleaning, and the finite certificate.  The formulation and its
correction history are summarized in \cite{BloomErdos796}.

The Mertens component follows the Abel-summation route in Goldmakher's
\emph{A quick proof of Mertens' theorem} \cite{GoldmakherMertens}, as
formalized in the \texttt{PrimeNumberTheoremAnd} library.  The accompanying
Lean development, including the explicit numerical certificates, is
available at
\url{https://github.com/rishigajjala/erdos-796-lean}.

Throughout the paper, all logarithms are natural.  The functions \(\pi(x)\)
and \(P^+(a)\) denote the prime-counting function and the largest prime
factor of \(a\), respectively, with \(P^+(1)=1\).  An integer is
\(z\)-smooth if all its prime factors are at most \(z\); thus \(1\) is
\(z\)-smooth.  Constants implicit in \(O(\cdot)\) and \(\ll\) are absolute
unless a dependence is displayed.

\section{Compatible cofactor families}\label{sec:model}

For finite sets \(U,V\subseteq\N\), define the ordered product convolution
\[
  r_{U,V}(m)
  =\#\{(u,v)\in U\times V:uv=m\}.
\]

\begin{definition}\label{def:compatible}
A sequence \(\Ucal=(U_j)_{j\geq1}\) is \emph{compatible} if
\[
  U_j\subseteq[j]
  \quad\text{and}\quad
  r_{U_i,U_j}(m)\leq2
  \qquad(i,j,m\geq1).
\]
\end{definition}

When \(i=j\), compatibility says that \(U_i\) has at most one unordered
factorization of any integer, with diagonal factorizations included.  Indeed,
two distinct off-diagonal factorizations give four ordered pairs, whereas a
diagonal and an off-diagonal factorization give three.

For \(j\geq1\), put
\begin{equation}\label{eq:Nj}
  N_j(n)=\#\left\{\sqrt n<q\leq n:q\text{ prime and }
      \floor{\frac nq}=j\right\},
\end{equation}
and define
\begin{equation}\label{eq:Gn}
  G(n)=\max_{\Ucal\ {\rm compatible}}
       \sum_{j\geq1}N_j(n)\abs{U_j}.
\end{equation}
Only \(j<\sqrt n\) have nonzero weight, so \eqref{eq:Gn} is a finite
optimization problem.

\begin{lemma}[The model lower bound]\label{lem:model-lower}
For every compatible family \(\Ucal\), the set
\begin{equation}\label{eq:lift}
  A_n(\Ucal)=
  \left\{qu:\sqrt n<q\leq n\text{ prime},\
     u\in U_{\floor{n/q}}\right\}
\end{equation}
is admissible for \(g_3(n)\).  Consequently \(G(n)\leq g_3(n)\).
\end{lemma}

\begin{proof}
If \(u\leq n/q<\sqrt n\), then \(q\) is the unique prime divisor of \(qu\)
exceeding \(\sqrt n\).  In particular, the parametrization in
\eqref{eq:lift} is injective.

Consider a product of two members of \(A_n(\Ucal)\).  If their large primes
are distinct, say \(q\) and \(Q\), every representation identifies which
factor contains \(q\) and which contains \(Q\).  It therefore injects into
\[
  U_{\floor{n/q}}\times U_{\floor{n/Q}},
\]
and compatibility permits at most two ordered cofactor pairs.  If both
large primes equal \(q\), representations correspond to unordered
off-diagonal factorizations inside one \(U_j\), of which there is at most
one.  All factor pairs counted here consist of distinct integers, as
required in the definition of \(g_3(n)\).
\end{proof}

The reverse inequality, up to the required error, is the structural heart
of the proof.

\begin{theorem}[Structural reduction]\label{thm:structural}
As \(n\to\infty\),
\[
  g_3(n)=G(n)+o\left(\frac n{\log n}\right).
\]
\end{theorem}

\section{A complete-box estimate}\label{sec:box}

Let \(K^{(3)}_{2,2,2}\) denote the complete tripartite \(3\)-uniform
hypergraph with two vertices in each part.

\begin{lemma}[Unbalanced complete-box bound]\label{lem:box}
If a tripartite \(3\)-uniform hypergraph \(H\), with positive part sizes
\(n_1,n_2,n_3\), contains no \(K^{(3)}_{2,2,2}\), then
\[
  e(H)\ll
  \frac{n_1n_2n_3}{\min(n_1,n_2,n_3)^{1/4}}.
\]
\end{lemma}

\begin{proof}
Write the part sizes as \(x\leq y\leq z\), and let \(d(u,v)\) be the
codegree of \((u,v)\in Y\times Z\) into the smallest part.  For each pair of
vertices in the smallest part, its common link in \(Y\times Z\) is
\(C_4\)-free.  Summing the K\H{o}v\'ari--S\'os--Tur\'an bound
\cite{KovariSosTuran1954} over those pairs gives
\[
  S:=\sum_{Y\times Z}\binom{d(u,v)}2\leq x^2z\sqrt y.
\]
Cauchy--Schwarz then gives
\[
  \frac{e(H)^2}{yz}
  \leq\sum_{Y\times Z}d(u,v)^2
  =e(H)+2S.
\]
Therefore
\[
  e(H)\leq yz+\sqrt2\,xy^{3/4}z
  \leq(1+\sqrt2)x^{3/4}yz,
\]
as required.
\end{proof}

\section{Structural reduction}\label{sec:structural}

Put \(L=\log n\) and
\begin{equation}\label{eq:parameters}
  Y=L^{24},\qquad Z=L,\qquad W=L^{1/8},\qquad R=Y^4=L^{96}.
\end{equation}
For large \(n\),
\[
  1<W<Z<Y<R<\sqrt n.
\]

We use one simple observation repeatedly.  Suppose triples
\((x,y,z)\) parametrize distinct elements \(xyz\in A\).  Their tripartite
hypergraph cannot contain \(K^{(3)}_{2,2,2}\): the eight elements in such a
cube would give the four complementary representations
\[
\begin{split}
 &(x_0y_0z_0)(x_1y_1z_1),\qquad
 (x_0y_0z_1)(x_1y_1z_0),\\
 &(x_0y_1z_0)(x_1y_0z_1),\qquad
 (x_0y_1z_1)(x_1y_0z_0)
\end{split}
\]
of the same integer.  Injectivity makes these four distinct-factor
representations.

\begin{lemma}[Pruning to two normal forms]\label{lem:pruning}
After deleting \(o(n/L)\) elements from an admissible \(A\subseteq[n]\),
every remaining element has exactly one of the forms
\[
\begin{array}{ll}
{\rm (S)}& qt,\quad q>\sqrt n\text{ prime},\quad
                 t<R,\quad P^+(t)\leq Z;\\[2mm]
{\rm (P)}& sqr,\quad s<W,\quad q>r>Z\text{ primes},\quad q>Y.
\end{array}
\]
All estimates are uniform in \(A\).
\end{lemma}

\begin{proof}
First discard elements having a factorization \(xyz\) with
\(x,y,z>Y\).  Choose one such factorization for each element and place its
coordinates in full dyadic ambient intervals.  By the observation above and
Lemma~\ref{lem:box}, each box contains \(O(nY^{-1/4})\) elements.  There are
\(O(L^3)\) boxes, so the discarded set has size
\[
  O(nL^3Y^{-1/4})=O(n/L^3)=o(n/L).
\]

A remaining \(Y\)-smooth integer is at most \(Y^6\).  Otherwise we could
greedily extract two divisors in \((Y,Y^2]\), leaving a third factor larger
than \(Y\).  Thus all remaining \(Y\)-smooth integers contribute at most
\[
  Y^6=L^{144}=o(n/L).
\]

For every other element put \(q=P^+(a)>Y\) and \(t=a/q\).  If \(t\) is
\(Z\)-smooth, then \(t<R\).  Indeed, if \(t\geq Y^4\), a minimal subproduct
\(d>Y\) satisfies \(d\leq YZ<Y^2\), and \(q,d,t/d\) are three factors
larger than \(Y\).  The cases \(q\leq\sqrt n\) contribute at most
\(R\pi(\sqrt n)=o(n/L)\), leaving form \({\rm (S)}\).

If \(t\) is not \(Z\)-smooth, put \(r=P^+(t)>Z\) and \(s=t/r\).  Then
\(q\geq r\) and \(s<R\).  To see the latter, suppose \(s\geq Y^4\).
If \(r>Y\), the factors \(q,r,s\) are all larger than \(Y\).  If
\(r\leq Y\), then \(s\) is \(Y\)-smooth; a minimal divisor
\(Y<d\leq Y^2\) makes \(q,d,rs/d\) three factors larger than \(Y\).
Both cases were already discarded.  The case \(q=r\) contributes at most
\(R\pi(\sqrt n)=o(n/L)\).

It remains to discard \(W\leq s<R\).  The triples \((s,q,r)\) are
canonical, since \(q=P^+(a)\), \(r=P^+(a/q)\), and \(q>r\); hence they
parametrize their products injectively.  In a dyadic box of scales
\(S,Q,P\), the three full ambient parts have upper orders
\[
  O(S),\qquad O(Q/\log Q),\qquad O(P/\log P),
\]
and the prime number theorem in dyadic intervals gives the corresponding
lower orders for the two prime parts.  Since \(S\geq W\), \(P\geq Z\), and
\(Q\geq Y\), each part has size \(\gg W\).  Lemma~\ref{lem:box} gives
\[
  e(S,Q,P)\ll
  \frac{SPQ}{\log P\log Q}\,W^{-1/4}.
\]
In every nonempty box \(Q\gg P\), so \(P\ll\sqrt{n/S}\) and therefore
\(\log(n/(SP))\geq\tfrac12\log(n/S)+O(1)\asymp L\).  For fixed \(S,P\),
summing the geometric \(Q\)-scales gives
\[
  \sum_Q e(S,Q,P)\ll\frac n{L\log P}\,W^{-1/4}.
\]
There are \(O(\log\log n)\) dyadic \(S\)-scales, and summing
\(1/\log P\) over the dyadic \(P\)-scales contributes another
\(O(\log\log n)\).  Hence
\[
  \sum_{S,Q,P}e(S,Q,P)
  \ll \frac nL\,W^{-1/4}(\log\log n)^2=o(n/L).
\]
The remaining elements are exactly those of forms \({\rm (S)}\) and
\({\rm (P)}\).
\end{proof}

\subsection{The semiprime form}

Let \(A_{\rm sp}\) be the elements of form \({\rm (P)}\), and for \(s<W\)
put
\[
  I_s=\{qr:sqr\in A_{\rm sp}\}.
\]
The canonical factorization is unique, so
\[
  |A_{\rm sp}|=\sum_{s<W}|I_s|.
\]

\begin{lemma}[Pairwise overlap]\label{lem:overlap}
For \(s\neq t<W\),
\[
  |I_s\cap I_t|
  \ll \frac{n}{L\sqrt{Z\log Z}}+\frac{\sqrt n}{L}.
\]
\end{lemma}

\begin{proof}
Use two labeled prime classes and orient every edge \(qr\in I_s\cap I_t\)
by \(q>r\).  The resulting simple bipartite graph is \(C_4\)-free.  A cycle
with edges \(q_ir_j\), \(i,j\in\{1,2\}\), would put the eight integers
\(sq_ir_j,tq_ir_j\) in \(A\), and the four complementary pairs would all
represent \(stq_1q_2r_1r_2\).  Since
\(s,t<W<Z<q_i,r_j\), all eight integers are distinct.  Moreover, all four
edges satisfy \(q_i>r_j\), so the two labeled vertex classes are numerically
disjoint.  Thus these are four genuine distinct-factor representations,
a contradiction.

Fix a dyadic scale \(P\) for the smaller prime.  The two vertex classes
have sizes \(O(P/\log P)\) and \(O(n/(PL))\).  The \(C_4\)-free estimate
therefore gives
\[
  O\left(\frac{n}{L\sqrt{P\log P}}+\frac{P}{\log P}\right)
\]
edges.  Summing over dyadic \(P\geq Z\) proves the lemma.
\end{proof}

Bonferroni's inequality and Lemma~\ref{lem:overlap} give
\[
\begin{aligned}
 |A_{\rm sp}|
 &\leq \left|\bigcup_{s<W}I_s\right|
   +O\left(W^2\left(
      \frac{n}{L\sqrt{Z\log Z}}+\frac{\sqrt n}{L}\right)\right)\\
 &=\left|\bigcup_{s<W}I_s\right|+o(n/L),
\end{aligned}
\]
because \(W^2=L^{1/4}\) and \(Z=L\).  The union is contained in the set of semiprimes
\(qr\leq n\) with \(q>Y\) and \(r>Z\).  Those with both primes at most
\(\sqrt n\) contribute \(O(n/L^2)\).  Otherwise the larger prime \(Q\)
exceeds \(\sqrt n\); if \(j=\lfloor n/Q\rfloor\), the other prime lies in
\((Z,j]\).  Hence
\begin{equation}\label{eq:Asp}
 |A_{\rm sp}|
 \leq \sum_{j\geq1}N_j(n)
       \bigl(\pi(j)-\pi(\min(j,Z))\bigr)+o(n/L).
\end{equation}

\subsection{The smooth form}

Let \(A_{\rm sm}\) be the elements of form \({\rm (S)}\), and define
\[
  T_q=\{t:qt\in A_{\rm sm}\}\qquad(q>\sqrt n).
\]
For distinct primes \(q,Q\),
\begin{equation}\label{eq:cross}
  r_{T_q,T_Q}(m)\leq2.
\end{equation}
Indeed, three ordered pairs \((u,v)\) with \(uv=m\) would give three
representations \((qu)(Qv)=qQm\).  These representations have unequal
factors and are distinct because \(q,Q>\sqrt n>R\).

An \emph{internal collision} is a pair of distinct unordered
\(2\)-multisets
\[
  \{a,b\}\neq\{c,d\},\qquad ab=cd,\qquad a,b,c,d\leq R.
\]
Diagonals such as \(\{a,a\}\) are allowed.  A fixed collision occurs in at
most one fiber.  If it occurred in both \(T_q\) and \(T_Q\), two
off-diagonal factorizations would give four ordered cross pairs, while a
diagonal together with an off-diagonal factorization would give three.
Both contradict \eqref{eq:cross}.

There are at most \(R^4\) possible collisions.  For each collision that
occurs, delete one tagged element \((q,x)\), meaning the copy of a support
value \(x\) in the unique fiber \(T_q\) containing that collision, and call
the cleaned fibers \(S_q\).  Then
\begin{equation}\label{eq:cleaned}
  \sum_q|T_q\setminus S_q|\leq R^4=o(n/L),
\end{equation}
and the whole family \((S_q)\), including each individual fiber, is
compatible.

For each \(j\) with \(N_j(n)>0\), choose \(S_j\) of maximum size among the
cleaned fibers in the bucket \(\{q:\lfloor n/q\rfloor=j\}\); if
\(N_j(n)=0\), put \(S_j=\varnothing\).  Choices from distinct nonempty
buckets come from distinct prime fibers, so compatibility is preserved, and
\[
  |A_{\rm sm}|
  \leq\sum_jN_j(n)|S_j|+o(n/L).
\]
Finally set
\[
  U_j=S_j\cup\{p:Z<p\leq j,\ p\text{ prime}\}.
\]
The two parts are disjoint.  The family \((U_j)\) is compatible because a
product contains zero, one, or two prime factors exceeding \(Z\); these
cases cannot collide, and compatibility or unique factorization gives at
most two orientations within each case.  Thus
\[
  |U_j|=|S_j|+\pi(j)-\pi(\min(j,Z)).
\]

Restoring the \(o(n/L)\) elements deleted in Lemma~\ref{lem:pruning}, and
combining this identity with \eqref{eq:Asp} and \eqref{eq:cleaned}, gives
\[
  |A|\leq\sum_jN_j(n)|U_j|+o(n/L)\leq G(n)+o(n/L).
\]
Together with Lemma~\ref{lem:model-lower}, this proves
Theorem~\ref{thm:structural}.

\section{The limit of the cofactor model}\label{sec:limit}

For a compatible family \(\Ucal=(U_j)\), put
\[
  w_j=\frac1{j(j+1)},\qquad
  e_j(\Ucal)=|U_j|-1-\pi(j).
\]

\begin{definition}\label{def:gamma}
Define
\[
  \Gamma=\sup_{\Ucal\ {\rm compatible}}
  \sum_{j\geq1}w_je_j(\Ucal),
\]
where a sum is assigned the value \(-\infty\) if its negative part diverges.
\end{definition}

Compatibility with \(i=j\) makes each \(U_j\) a multiplicative Sidon set
with diagonals allowed.  The classical bound of Erd\H{o}s
\cite{Erdos1938Sequences}, also recorded in
\cite[p.~252]{Erdos1964Multiplicative}, is
\[
  |U_j|\leq\pi(j)+O(j^{3/4})
\]
and therefore gives
\begin{equation}\label{eq:positive}
  e_j^+\ll j^{3/4},
  \qquad \sum_{j\geq1}w_je_j^+<\infty,
\end{equation}
uniformly over compatible families.  The family consisting of \(1\) and all
primes has \(e_j=0\), so \(0\leq\Gamma<\infty\).

\begin{proposition}[An explicit upper bound]\label{prop:explicit-upper}
One has
\[
  \Gamma
  \leq \zeta\left(\frac76\right)
      +\frac32\zeta\left(\frac43\right)
  <13.
\]
Consequently
\[
  1+M+\Gamma<15.
\]
\end{proposition}

\begin{proof}
We use the standard factor-graph argument, as in
\cite[Section~3.2]{LiuPach2019}.  Let \(U\subseteq[N]\) be
multiplicative Sidon.

First consider those \(a\in U\) having a prime divisor
\(p>N^{2/3}\).  Such a prime is unique, and we write \(a=pb\), where
\(b<N^{1/3}\).  Form the bipartite graph whose edges are these pairs
\((p,b)\).  Distinct edges represent distinct integers, and the graph is
\(C_4\)-free: a four-cycle would give four distinct elements satisfying
\[
  (pb)(qc)=(pc)(qb).
\]
If \(d_p\) is the degree of \(p\), then
\[
  \sum_p\binom{d_p}{2}
  \leq\binom{\floor{N^{1/3}}}{2}.
\]
Since \(d_p-1\leq\binom{d_p}{2}\), this part of \(U\) has at most
\[
  \pi(N)-\pi(N^{2/3})+\frac12N^{2/3}
\]
elements.

Every remaining \(a\leq N\) has a factorization
\[
  a=uv,\qquad v\leq u\leq N^{2/3}.
\]
Indeed, this is immediate when \(a\leq N^{2/3}\).  Otherwise, if \(a\)
has a prime divisor \(N^{1/3}\leq p\leq N^{2/3}\), take the factors \(p\) and
\(a/p\).  If all prime divisors are smaller than \(N^{1/3}\), multiply
prime factors until the product first exceeds \(N^{1/3}\); that product
and its complementary factor are both at most \(N^{2/3}\).  After
interchanging the factors if necessary, also
\(v\leq\sqrt a\leq N^{1/2}\).

Fix one such factorization for each remaining element and form the
corresponding bipartite graph.  Its two vertex classes have sizes at most
\(\floor{N^{2/3}}\) and \(\floor{N^{1/2}}\).  Distinct edges represent distinct
integers, so a four-cycle would again contradict the multiplicative
Sidon property.  The standard \(C_4\)-free estimate
\cite{KovariSosTuran1954} gives
\[
  e\leq N^{1/2}\sqrt{N^{2/3}}+N^{2/3}
    =N^{5/6}+N^{2/3}.
\]
Combining the two parts gives
\[
  |U|-1-\pi(N)
  \leq N^{5/6}+\frac32N^{2/3}.
\]

Applying this to every fiber \(U_j\), and using \(w_j<j^{-2}\), gives
\[
  \Gamma
  \leq\zeta\left(\frac76\right)
       +\frac32\zeta\left(\frac43\right).
\]
For \(s>1\),
\[
  \zeta(s)<1+\frac1{s-1},
\]
so
\[
  \Gamma<7+\frac32\cdot4=13.
\]
Finally, the explicit cutoff estimate in the accompanying Lean development
gives \(M<933/1000<1\), and therefore
\[
  1+M+\Gamma<15.
\]
\end{proof}

\begin{lemma}[Canonical extension]\label{lem:extension}
Every compatible prefix \(U_1,\ldots,U_J\) extends compatibly by
\[
  U_j=U_J\cup\{p:J<p\leq j,\ p\text{ prime}\}\qquad(j>J).
\]
For this extension, \(e_j=e_J\) for \(j>J\).
\end{lemma}

\begin{proof}
Old elements are at most \(J\), while every new prime exceeds \(J\).
For two fibers, products with no new prime are controlled by
\(r_{U_i,U_J}\) or \(r_{U_J,U_J}\), hence by the original prefix.  With one
new prime, unique factorization determines that prime and the old factor,
leaving at most two orientations; with two new primes it again leaves at
most two.  Finally,
\[
 |U_j|=|U_J|+\pi(j)-\pi(J),
\]
so \(e_j=e_J\).
\end{proof}

Since \(\sum_{j>J}w_j=1/(J+1)\), the extension lemma implies
\begin{equation}\label{eq:prefix}
  \sum_{j\leq J}w_je_j
  \leq\Gamma+\frac{1+\pi(J)}{J+1}
  =\Gamma+O(1/\log J).
\end{equation}
Set
\[
  B(n)=\sum_{j\geq1}N_j(n)(1+\pi(j)).
\]

\begin{proposition}\label{prop:model-limit}
As \(n\to\infty\),
\[
  G(n)=B(n)+(\Gamma+o(1))\frac n{\log n}.
\]
\end{proposition}

\begin{proof}
For fixed \(j\), the prime number theorem gives
\begin{equation}\label{eq:weight-limit}
  \frac{\log n}{n}N_j(n)\longrightarrow
  \frac1j-\frac1{j+1}=w_j.
\end{equation}
The positive tails are uniform.  In a dyadic block \(K\leq j<2K\),
\eqref{eq:positive} and Chebyshev's bound give
\[
  \sum_{K\leq j<2K}N_j(n)e_j^+
  \ll K^{3/4}\pi(n/K)
  \ll\frac n{\log n}K^{-1/4}.
\]
Summing the blocks yields
\begin{equation}\label{eq:tail}
  \frac{\log n}{n}\sum_{j>J}N_j(n)e_j^+
  \ll J^{-1/4}.
\end{equation}

Apply this to a family attaining \(G(n)\).  For fixed \(J\),
\eqref{eq:weight-limit} is uniform over its prefix because
\(|e_j|\leq2j+1\).  Hence \eqref{eq:prefix} and \eqref{eq:tail} give
\[
  \limsup_{n\to\infty}\frac{\log n}{n}(G(n)-B(n))
  \leq\Gamma+O(1/\log J)+O(J^{-1/4}).
\]
Let \(J\to\infty\).

Conversely, choose a compatible family whose defining sum exceeds
\(\Gamma-\varepsilon\).  Its defining series is absolutely convergent: the
positive part is finite by \eqref{eq:positive}, and the negative part is
finite by Definition~\ref{def:gamma}.  Moreover, the trivial lower bound
\(e_J\geq-1-\pi(J)\),
\[
  \frac{|e_J|}{J+1}\ll J^{-1/4}+\frac1{\log J}.
\]
Choose \(J\) so that
\[
  \left|\sum_{j>J}w_je_j\right|<\varepsilon
  \qquad\text{and}\qquad
  \frac{|e_J|}{J+1}<\varepsilon.
\]
The canonical extension of this prefix then has weighted value
\[
  \sum_{j\leq J}w_je_j+\frac{e_J}{J+1}
  >\Gamma-3\varepsilon.
\]
For this fixed extension,
\[
  \sum_{j>J}N_j(n)
  =\pi\left(\frac n{J+1}\right)-\pi(\sqrt n)
\]
once \(n>(J+1)^2\).  Together with \eqref{eq:weight-limit}, the prime number
theorem gives
\[
  \frac{\log n}{n}\sum_jN_j(n)e_j
  \longrightarrow\sum_jw_je_j.
\]
Thus the lower limit is at least
\(\Gamma-3\varepsilon\).  Letting \(\varepsilon\downarrow0\) proves the
proposition.
\end{proof}

It remains to evaluate \(B(n)\).  Let \(\pi_2(n)\) count products
\(pq\leq n\) with \(p\leq q\) prime, including squares.  Direct counting
gives
\begin{align}
  \sum_jN_j(n)&=\pi(n)-\pi(\sqrt n),\label{eq:prime-count}\\
  \sum_jN_j(n)\pi(j)
   &=\pi_2(n)-\frac{\pi(\sqrt n)(\pi(\sqrt n)+1)}2.
   \label{eq:semiprime-count}
\end{align}
The semiprime asymptotic
\cite[Theorem~1, equation~(17)]{CrisanErban2021} is
\[
  \pi_2(n)
  =\frac{n\log\log n}{\log n}
   +(M+o(1))\frac n{\log n}.
\]
Consequently,
\begin{equation}\label{eq:B-final}
  B(n)=\frac{n\log\log n}{\log n}
       +(1+M+o(1))\frac n{\log n}.
\end{equation}
Theorem~\ref{thm:structural}, Proposition~\ref{prop:model-limit}, and
\eqref{eq:B-final} prove the asymptotic formula in
Theorem~\ref{thm:main}.  The bound \(\Gamma\geq4/15\) is proved next.

\section{An explicit lower bound}\label{sec:certificate}

Define the following seven fibers:
\[
\begin{array}{lll}
  \mathsf A=\{1\},&
  \mathsf B=\{1,2\},&
  \mathsf C=\{1,2,3\},\\[2pt]
  \mathsf D=\{1,3,4,5\},&
  \mathsf E=\{1,3,4,5,6\},&
  \mathsf F=\{1,3,4,5,6,7\},\\[2pt]
  &&\mathsf G=\{2,3,5,7,8,9,10\}.
\end{array}
\]
Assign the fibers
\[
  \mathsf A,\mathsf B,\mathsf C,\mathsf C,\mathsf D,
  \mathsf E,\mathsf F,\mathsf F,\mathsf F,\mathsf G
\]
to \(j=1,\ldots,10\).  The complete multiplication certificate in
Appendix~\ref{app:certificate} proves that this prefix is compatible.  Its excess
sequence is
\[
  (e_1,\ldots,e_{10})=(0,0,0,0,0,1,1,1,1,2).
\]
The weighted prefix value is
\[
  \frac1{42}+\frac1{56}+\frac1{72}+\frac1{90}+\frac1{55}
  =\frac{14}{165}.
\]
The canonical extension from \(\mathsf G\) has \(e_j=2\) for \(j>10\),
and contributes \(2/11\).  Therefore
\begin{equation}\label{eq:four-fifteenths}
  \Gamma\geq\frac{14}{165}+\frac2{11}
  =\frac4{15}.
\end{equation}
This completes the proof of Theorem~\ref{thm:main}.

\appendix

\section{The multiplication certificate}\label{app:certificate}

For two fibers \(\mathsf X,\mathsf Y\), let
\[
  D(\mathsf X,\mathsf Y)
  =\{m:r_{\mathsf X,\mathsf Y}(m)=2\}.
\]
For each unordered pair of fiber types, the following table lists all
products occurring more than once.  Every listed product has convolution
coefficient exactly two, and direct multiplication shows that every
remaining product occurs at most once.  Thus all \(28\) fiber-pair
convolutions satisfy the compatibility condition.

\begingroup
\small
\renewcommand{\arraystretch}{1.12}
\noindent
\begin{tabularx}{\textwidth}{
  >{\raggedright\arraybackslash}p{0.18\textwidth}
  >{\raggedright\arraybackslash}X}
\toprule
Pair & Products with convolution coefficient exactly \(2\)\\
\midrule
\(\mathsf{AA},\mathsf{AB},\mathsf{AC},\)\newline
\(\mathsf{AD},\mathsf{AE},\mathsf{AF},\mathsf{AG}\) & --\\
\(\mathsf{BB}\) & \(2\)\\
\(\mathsf{BC}\) & \(2\)\\
\(\mathsf{BD}\) & --\\
\(\mathsf{BE}\) & \(6\)\\
\(\mathsf{BF}\) & \(6\)\\
\(\mathsf{BG}\) & \(10\)\\
\(\mathsf{CC}\) & \(2,3,6\)\\
\(\mathsf{CD}\) & \(3\)\\
\(\mathsf{CE}\), \(\mathsf{CF}\) & \(3,6,12\)\\
\(\mathsf{CG}\) & \(6,9,10\)\\
\(\mathsf{DD}\), \(\mathsf{DE}\), \(\mathsf{DF}\)
  & \(3,4,5,12,15,20\)\\
\(\mathsf{DG}\) & \(8,9,10,15,40\)\\
\(\mathsf{EE}\), \(\mathsf{EF}\)
  & \(3,4,5,6,12,15,18,20,24,30\)\\
\(\mathsf{EG}\) & \(8,9,10,12,15,30,40\)\\
\(\mathsf{FF}\)
  & \(3,4,5,6,7,12,15,18,20,21,24,28,30,35,42\)\\
\(\mathsf{FG}\)
  & \(8,9,10,12,15,21,30,35,40\)\\
\(\mathsf{GG}\)
  & \(6,10,14,15,16,18,20,21,24,27,30,35,40,\)\newline
    \(45,50,56,63,70,72,80,90\)\\
\bottomrule
\end{tabularx}
\endgroup

% The default AMS shorthand for repeated authors is a horizontal rule.
% Here the only repeated author is Erd\H{o}s, so print his name in full.
\providecommand{\bysame}{Paul Erd\H{o}s}
\begingroup
\footnotesize
\bibliographystyle{amsplain}
\bibliography{references}
\endgroup

\end{document}
